% ============================================================================
\chapter{Troubleshooting}
\label{ch:troubleshooting}
% ============================================================================

\section{Common problems}

\begin{center}
\begin{tabularx}{\textwidth}{@{}p{0.34\textwidth}X@{}}
  \toprule
  \textbf{Symptom} & \textbf{Cause and remedy} \\
  \midrule
  Status: \emph{Connection to \dots\ failed: Permission denied} &
    The user lacks serial-port rights. \code{sudo usermod -aG dialout <user>},
    then log out/in (or restart the service). \\
  \addlinespace
  Status: \emph{Connection \dots\ failed: No such file or directory} &
    Wrong \code{port} in \code{config.toml}, or the USB cable is unplugged.
    Use \textbf{List com} / \code{ls /dev/serial/by-id/} to identify the device. \\
  \addlinespace
  Status: \emph{Device on \dots\ did not respond correctly to IsDelta} &
    A serial device was found but it is not (or not yet) the robot: wrong port,
    wrong baud rate, or robot still booting. Verify \code{baud\_rate = 115200},
    power-cycle the robot, restart the application. \\
  \addlinespace
  Status: \emph{Timeout waiting for 'ok' from robot} &
    The robot stopped answering mid-command: cable disturbed, firmware halted, or
    mechanical fault. Check the machine, then restart the application and re-home
    before continuing. \\
  \addlinespace
  Status: \emph{Movement out of safety limits} &
    The requested jog would leave the configured box. If the limits are too tight
    for your tooling, adjust \code{[robot]} in \code{config.toml}
    (Section~\ref{sec:limits}). \\
  \addlinespace
  Application does not start on the Pi (blank screen) &
    Check \code{journalctl -u deltax2}. Typical causes: missing/invalid
    \code{config.toml} in the \code{WorkingDirectory}, or the binary was built
    without the \code{linuxkms} backend feature while
    \code{SLINT\_BACKEND=linuxkms} is set (Section~\ref{sec:rendering}). \\
  \addlinespace
  Touch input misaligned or rotated &
    Configure \code{LIBINPUT\_CALIBRATION\_MATRIX} via a udev rule
    (Section~\ref{sec:touchcal}). \\
  \bottomrule
\end{tabularx}
\end{center}

\section{Diagnostics over SSH}

\begin{lstlisting}[language=bash]
$ journalctl -u deltax2 -f          # live application log
$ ls -l /dev/serial/by-id/          # stable names of USB serial adapters
$ sudo systemctl restart deltax2    # restart the UI after config changes
\end{lstlisting}

The log contains every serial command sent (\code{Serial write: \dots}), all status
messages shown in the UI, and the output of \textbf{List com}.
